\section{Implementation-independent computational policy}
\label{sec:policy}
The preceding bounds suggest a reproducible decision procedure,
but not a universal numerical winner. In practice the expensive
operations must be attributed to the correct level of reuse:
a factorization is specific to $(n,d,\lambda)$, whereas
a smoothing-grid inverse is specific only to $(n,d,s)$.
A fitted curve is specific to the observations and to the
reference-difference mean. A scalar selection criterion
depends on the already-computed residual and diagnostic
summaries, whereas its numerical definition is unchanged
by a reparameterization of the candidate grid.

\subsection{Cost-aware organization of a complete search}
For each fixed $(n,d)$ the following abstract procedure
avoids unnecessary repetitions. It is pseudocode for
organizing computations, not a claim of a newly implemented
universal dispatcher.
\begin{enumerate}
\item Fix the requested criterion set $\mathcal C$,
the target-smoothness grid $\mathcal S$, numerical tolerances,
and fixed or explicitly conditional plug-in means.
Determine whether any requested criterion requires
per-observation leverages (ordinary CV) or a
Gaussian marginal determinant.
\item Choose a mathematically admissible smoothness
mapping method. Invert each interior $s_j$ \emph{once};
handle $s=0$ and $s=1$ by their respective analytical
limits. Record a bracketed inversion residual or a
rigorous stopping condition, not merely an iteration count.
\item Choose the candidate-score family subject to the
memory budget and required outputs. For the spectral
family, prepare the basis once and the projected
observations once per series. For the banded family,
factor each $A_{\lambda_j}$ once and obtain its
selected inverse or other needed diagnostics once.
\item For each $j=1,\ldots,M$, reuse the same penalty
and factorization across every observed series and
every criterion that permits such sharing. Calculate
all required score terms, flag undefined scores and
numerical failures, and update the best valid
candidate for each series--criterion pair.
\item Record the decomposition, grid mapping,
criterion evaluation, memory and accuracy costs
\emph{separately}; report the complete elapsed time
only for workloads with all required results.
\end{enumerate}

If no score requires the full fitted curve, computing
it for every candidate is unnecessary for the spectral
family. Conversely, ordinary conditional CV requires
leverage diagonals and residuals, so a fit-only solver
time is not a full CV-selection time. Marginal
Gaussian likelihood also requires its own variance
and determinant assumptions; the
least-squares objective should not be substituted
for the likelihood merely because both involve
$D_d^\top D_d$.

\subsection{Rival algorithmic explanations}
One explanation for good observed performance is
\emph{arithmetic sparsity}: the $d$-banded positive
definite system admits a direct solve and its
selected-inverse diagnostic in nearly linear
arithmetic for fixed $d$. This is the strong
case for banded Cholesky when $n$ is large.

A rival explanation is \emph{amortization}:
dense spectral preprocessing is expensive,
but the same singular basis can serve many
smoothed-series and scalar-criterion evaluations.
The second explanation is especially relevant
for large $M$ and scalar GCV/AICc/BIC/ML scoring.
Neither explanation dominates without the
actual workload and implementation constants,
and both can be affected by numerical conditioning.
The synthesis is the conditional admissible
frontier of Eq.~\eqref{eq:theory_frontier}, rather
than a globally preferred factorization.

The importance of exploiting Whittaker--Henderson
structure was demonstrated long before this paper:
\citet{weinert2007} developed efficient fitted
trends and GCV computations, and \citet{biessy2026}
examined banded methods and reduced-rank
parameterizations in a modern statistical framework.
The sparse selected-inverse principle is likewise
classical \citep{erisman1975}. Our contribution
is an explicit \emph{combined workload model}
that keeps normalized-smoothness inversion,
criterion diagnostics, reuse, endpoints,
numerical admissibility and memory on the
same accounting basis. It is not a claim
to have invented Cholesky, selected inversion,
the Gaussian variance model or the original
smoothness index.

\subsection{Limits of the theoretical results}
Big-$O$ bounds omit constants, BLAS scheduling,
cache locality, dense-SVD numerical tolerances,
Python dispatch, allocation, garbage collection
and CPU contention. The threshold in
Eq.~\eqref{eq:theory_threshold} can be fitted
only from controlled measurements; using a
few timings collected under different
concurrency conditions would be misleading.
Furthermore, the numerical advantage of
shift-free spectral reconstruction at
large $d$ or extreme $\lambda$ is a separate
accuracy concern, not an asymptotic runtime
guarantee.

The statements concern a finite, equally spaced
univariate series and a conventional forward
difference penalty. Irregularly spaced
observations, missing-data weights, nonquadratic
losses, spatial lattices and higher-dimensional
penalties can change the sparsity pattern
and invalidate a naive reuse of $O(nd^2)$.
Although the fitted trend is commonly described
as a time-series trend, no in-sample conditional
CV expression presented here establishes
out-of-sample forecasting accuracy.

The endpoints also require care:
$S(0)=0$ returns the observations, but
$\mathrm{CV}$ and information criteria can be
undefined or divergent at exact interpolation;
$S=1$ means the constrained polynomial limit,
not the inversion of $A_\lambda$ at an
artificially enormous finite penalty.
A user-provided $\mu$ is fixed when the
derivatives and conditional EDF are computed.
If $\mu$ is estimated from $y$, this
data dependence cannot simply be suppressed
when deriving the exact leave-one-out score.

\section{Conclusions}\label{sec:theory_conclusions}
Smoothness-indexed finite-difference PLS
admits two complementary computational
representations. Its reusable spectrum
supplies a monotone and data-independent
map between normalized smoothness and
the penalty; its banded normal equations
supply a memory-efficient way to obtain
fitted curves and exact conditional
criterion diagnostics. We derived a
complete banded expression for the
reduced-coordinate Gaussian determinant
and profiled quadratic form, established
amortized work and storage bounds for
repeated candidates, criteria and series,
and expressed the implementation decision
as an admissible, workload-specific
cost frontier.

These deductions do not determine a universal fastest method
in milliseconds. Section~\ref{sec:matched-example} reports
one matched, single-worker BIC implementation example;
the observed spectral advantage for a $10^5$-candidate grid
is descriptive at $n=64$ and $d=1$, not a claim about CV
or arbitrary sample sizes. The NGSIM application
in Section~\ref{sec:jerk-illustration} likewise demonstrates
the meaning of the estimated third difference and a
separately appended forecasting heuristic, not an
intrinsic PLS predictive distribution or measured true jerk.
A fully powered multi-criterion runtime frontier and
independent physical forecasting evaluation remain
appropriate subjects for separate papers.

\section*{Data and code availability}
The full mathematical derivations and
reproducibility notes are maintained at
\url{https://github.com/heri-espino/trendestimation}.
An archived source revision must be
cited at submission. The mathematical identities require no
experimental calibration. The small matched
benchmark illustration uses 48 frozen
Windows observations and their original source revision,
copied verbatim under
\texttt{paper/csda-computational-costs/data/}.
The jerk example reuses real FHWA US-101 longitudinal
measurements and the predeclared training-only selection
from \texttt{paper/ngsim-jerk-case-study/};
the recorded positions and CC BY-SA 3.0 attribution
are retained in that source manuscript.
The four new AwesomeFigures sources, editable
PGF/PDF/SVG outputs, figure-generation script,
checksums, and full archived benchmark provenance
are preserved in this repository.

\section*{Declarations}
Author order, correspondence,
funding, potential conflicts and
contributor statements remain subject
to confirmation by both authors
before journal submission.
